%============ abstract.tex ===================
\chapter*{Abstract}
\addcontentsline{toc}{chapter}{Abstract}

Parkinson's Disease impairs the neuromuscular control of speech,
producing measurable changes in the voice that can appear early in the
disease. This project develops and evaluates a complete, reproducible
research prototype that analyses raw continuous-speech recordings for
Parkinson's-related acoustic patterns using classical, explainable
machine learning, with particular emphasis on honest validation.

From each recording, after standardized preprocessing and segmentation
into 10-second chunks, 74 physiologically motivated acoustic features
are extracted --- fundamental-frequency statistics, jitter, shimmer,
harmonics-to-noise ratio, smoothed cepstral peak prominence, pause
statistics, and mel-frequency cepstral coefficients with their
dynamics --- and averaged per recording. Classifiers were compared on
the MDVR-KCL corpus (37 subjects) under strictly subject-independent
cross-validation, with model selection governed by rules declared
before experimentation. The adopted Random Forest achieved a
subject-level balanced accuracy of $0.822 \pm 0.032$, sensitivity of
0.771, specificity of 0.873, and an area under the receiver operating
characteristic curve of 0.864, with pitch range and spectral
variability as its leading features. A controlled comparison showed
that a deliberately leaky chunk-level split inflates the same
pipeline's result from 0.775 to 0.909, quantifying why
subject-independent validation is essential. External validation on an
independent Italian corpus (61 speakers), with the model frozen and
recording bandwidth harmonized to neutralize a discovered sample-rate
confound, yielded a speaker-level area under the curve of 0.701 and
balanced accuracy of 0.629 across language and recording conditions.
The accompanying screening prototype reports scores between 0.35 and
0.65 as inconclusive, attaches recording-condition warnings, and
frames every output as a non-diagnostic research indication requiring
professional clinical evaluation.
\thispagestyle{plain}
